\section{Related Work}

\subsection{Trust and reputation in multi-agent systems}

The dominant paradigm for handling source differences is trust and
reputation, and it is almost entirely about reliability. TRAVOS~\cite{teacy2006}
estimates the probability that a source provides accurate information;
the Beta reputation system maintains Bayesian posteriors over
trustworthiness; FIRE, Regret, and BLADE extend the idea to witness
reputation and roles. Every one of these models \emph{how often}
a source is right, a scalar reliability parameter. That is the gap we work
in: a source can be maximally reliable under every one of these
frameworks and still be unreadable, because reliability and convention are
different things.

\subsection{Opinion dynamics and social learning}

{\sloppy
Classical opinion dynamics, DeGroot~\cite{degroot1974} and
Friedkin--Johnsen~\cite{friedkin1999}, aggregate through nonnegative
influence weights: exactly the class our Theorem~\ref{thm:nonneg} rules
out for inverted conventions. Bayesian social learning~\cite{acemoglu2011}
asks when networks converge to the truth, but takes the mapping from
state to signal as fixed and common knowledge. We study what happens when
that mapping is private, heterogeneous, and has to be inferred.
\par}

\subsection{Truth discovery}

Truth discovery methods~\cite{dong2009,yin2008,pasternack2011} jointly
estimate latent truth and source reliability from conflicting claims. Each
source still gets a scalar weight, and the weights still determine the
aggregation. Convention heterogeneity breaks this parameterisation from
the inside: the correct ``weight'' of a source is not a scalar at all but
a context-dependent decoding of the message. Truth discovery asks who to
trust. We are asking what they meant.

\subsection{Robust and adversarial aggregation}

Byzantine-tolerant aggregation~\cite{blanchard2017} and defences against
adversarial federated learning~\cite{yin2018} handle sources that actively
try to corrupt the outcome. Our setting is different: nobody is attacking
anyone. The sources are honest and often highly reliable; they simply
speak in different conventions, so there is no attacker to detect. What
remains is a decoding problem, and reliability scores cannot see it.

\subsection{Information theory and pragmatics}

The framing grew out of partial information
decomposition~\cite{williams2010,bertschinger2014}: does source identity carry information beyond message
content? What survived is decision-theoretic -- our theorems target
decision rules directly, since statistically real information
($I(X;S\mid M) > 0$) need not imply decision value. The intuition is old
in pragmatics: Gricean implicature~\cite{grice1975} treats speaker
identity and context as constitutive of meaning; our conventions are a
minimal formal version of that insight.

\subsection{Multi-agent LLM systems}

Recent work on LLM agents has explored debate protocols~\cite{du2023debate,
liang2023debate}, round-table deliberation~\cite{chen2024reconcile},
multi-agent evaluation~\cite{chan2024chateval}, and layered aggregation of
proposals~\cite{wang2024moa}. These frameworks model source differences
through confidence scores, voting, or an aggregator model's judgement, and
empirical studies find debate gains sensitive to protocol design and not
always robust over simpler baselines~\cite{smit2024mad}. What they
generally do not model is the possibility that agents encode readings
through source-specific, context-dependent conventions. Our contribution
relative to this literature is the impossibility results:
Theorems~\ref{thm:nonneg} and~\ref{thm:scalar} show that when the mapping
from readings to messages itself varies by source and context, the thing
to learn is a decoder, not a weight.
